\section*{Hoofdstuk \ref{ch:b3:mantle-convection} --- Mantelconvectie en plaatdrijvende krachten}

\begin{solution}{exo:b3:mantle-convection:1}
$\mathrm{Ra} = 3500\times9.9\times\num{3e-5}\times1000\times(\num{6.7e5})^3/
(10^{-6}\times\num{5e20}) = \num{6.25e5}$, ongeveer 950 keer het \omterm{def:b3:mantle-convection:critical}{kritisch
Rayleighgetal} 657.5: de laag convecteert krachtig.
\end{solution}

\begin{solution}{exo:b3:mantle-convection:2}
$\tau_M = 10^{21}/\num{7.7e10} = \qty{1.30e10}{s}$, ongeveer \qty{412}{yr}. Het
getij, dat uren duurt, treft een elastische mantel; het smelten van een ijskap,
over duizenden jaren, en het zinken van een slab, over miljoenen, treffen een
viskeuze mantel.
\end{solution}

\begin{solution}{exo:b3:mantle-convection:3}
$\tau = 4\pi\times10^{21}/(3300\times9.8\lambda)$: \qty{3.89e11}{s} =
\qty{12.3}{kyr} voor \qty{1000}{km}, en \qty{4.1}{kyr} voor \qty{3000}{km}. De
bredere belasting herstelt sneller.
\end{solution}

\begin{solution}{exo:b3:mantle-convection:4}
$g\rho_m\alpha(T_m - T_0)\kappa t = 9.8\times3300\times\num{3.1e-5}\times1450
\times\num{0.8e-6}\times50\times\num{3.156e13} = \qty{1.835e12}{N.m^{-1}}$; de
factor tussen haakjes is $1 + 2\times3300\times\num{3.1e-5}\times1450/(\pi\times2270) = 1.0416$,
zodat $F_{rp} = \qty{1.91e12}{N.m^{-1}}$, 24 keer kleiner dan de slabtrek.
\end{solution}

\begin{solution}{exo:b3:mantle-convection:5}
(a) $v = 3300\times9.8\times\num{3e-5}\times300\times(\num{e4})^2/10^{21} =
\qty{2.91e-11}{m.s^{-1}}$; de stijgtijd over \qty{10}{km} is $d/v =
\qty{3.44e14}{s}$, \qty{10.9}{Myr}, de diffusietijd $d^2/\kappa =
\qty{e14}{s}$, \qty{3.2}{Myr}. Hun verhouding is $(d^2/\kappa)/(d/v) =
vd/\kappa = \rho g\alpha\Delta Td^3/(\kappa\eta) = \mathrm{Ra} = 0.29$. (b) Het
koelt eerst af: kleine blobs convecteren niet. (c) $\mathrm{Ra} \propto d^3$, dus $d =
10\times(657.5/0.291)^{1/3} = \qty{131}{km}$.
\end{solution}

\begin{solution}{exo:b3:mantle-convection:6}
$2\mathrm{Ra}_c = 1315$ en $\mathrm{Ra}/2\mathrm{Ra}_c = 7604$, waarvan de
derdemachtswortel is 19.66. $\delta = 2891/19.66 = \qty{147}{km}$, $\mathrm{Nu} = 19.66/2 =
9.8$, $u = (10^{-6}/\num{2.891e6})\times19.66^2 = \qty{1.34e-10}{m.s^{-1}} =
\qty{0.42}{cm/yr}$, en $q = 9.8\times3.1\times1450/\num{2.891e6} =
\qty{15}{mW.m^{-2}}$.
\end{solution}

\begin{solution}{exo:b3:mantle-convection:7}
$u = 0.09/\num{3.156e7} = \qty{2.85e-9}{m.s^{-1}}$ en $B = 2270\times\num{8e8}
\times\num{2.85e-9} = \qty{5.18e3}{kg.s^{-1}} = \qty{5.2}{Mg/s}$, ongeveer
drie vijfde van die van Hawaï. De warmtestroom is $1170\times5178/\num{3.1e-5} =
\qty{1.95e11}{W}$, \qty{0.2}{TW}, of \qty{0.42}{\percent} van het
warmteverlies van de Aarde.
\end{solution}

\begin{solution}{exo:b3:mantle-convection:8}
$v = 2\times30\times9.8\times(\num{4e5})^2/(9\times10^{21}) =
\qty{1.05e-8}{m.s^{-1}} = \qty{0.33}{m/yr}$; \qty{2000}{km} nemen
$\num{2e6}/\num{1.05e-8} = \qty{1.9e14}{s}$, ongeveer \qty{6}{Myr}.
\end{solution}

\begin{solution}{exo:b3:mantle-convection:9}
$w = 300e^{-10/4.4} = \qty{30.9}{m}$ en $\dot w = w/\tau = 30.9/4400 =
\qty{7.0}{mm/yr}$, twee derde van de nabij Umeå gemeten snelheid: een uniforme
halfruimte met deze $\tau$ ligt in het juiste bereik, maar het echte herstel is
sneller of de indeuking dieper.
\end{solution}

\begin{solution}{exo:b3:mantle-convection:10}
(a) Met $w = 2\rho_m\alpha\Delta T\sqrt{\kappa t/\pi}/(\rho_m - \rho_w)$,
\[
  \tfrac12(\rho_m - \rho_w)gw^2 = \frac{2g\rho_m^2\alpha^2\Delta T^2\kappa t}
  {\pi(\rho_m - \rho_w)} = g\rho_m\alpha\Delta T\kappa t\times
  \frac{2\rho_m\alpha\Delta T}{\pi(\rho_m - \rho_w)} .
\] (b) Door partiële integratie,
$\int_0^\infty x\operatorname{erfc}x\,\mathrm dx = [\tfrac{x^2}2\operatorname{erfc}
x]_0^\infty + \int_0^\infty\tfrac{x^2}2\,\tfrac2{\sqrt\pi}e^{-x^2}\,\mathrm dx =
0 + \tfrac1{\sqrt\pi}\cdot\tfrac{\sqrt\pi}4 = \tfrac14$. (c) Op \qty{100}{Myr},
$\rho_mg\alpha\int z\theta = \qty{3.67e12}{N.m^{-1}}$ en, met $w =
\qty{3705}{m}$, $\tfrac12(\rho_m - \rho_w)gw^2 = \qty{1.53e11}{N.m^{-1}}$: de
waterterm is slechts \qty{4}{\percent} van de kracht.
\end{solution}

\begin{solution}{exo:b3:mantle-convection:11}
$\ln f = 3\ln(\pi^2 + k^2) - 2\ln k$, waarvan de afgeleide $6k/(\pi^2 + k^2) - 2/k$
verdwijnt wanneer $3k^2 = \pi^2 + k^2$, $k^2 = \pi^2/2$. Daar $f = (3\pi^2/2)^3/
(\pi^2/2) = 27\pi^4/4$. De golflengte $2\pi/k = 2\sqrt2\,d = 2.83\,d$ omvat een
paar tegengesteld draaiende cellen, elk $\sqrt2\,d = 1.41\,d$ breed.
\end{solution}

\begin{solution}{exo:b3:mantle-convection:12}
(a) $u = \qty{1.90e-9}{m.s^{-1}}$, $\eta u/h = 10^{20}\times\num{1.90e-9}/
\num{1.5e5} = \qty{1.27}{MPa}$, en over \qty{6000}{km}, \qty{7.6e12}{N.m^{-1}}.
(b) Twee keer de rugduw en een zesde van de slabtrek: de \omterm{def:b3:mantle-convection:basal-drag}{basale weerstand} kan
alleen de rugduw weerstaan, maar niet de slab. (c) $\eta uL/h = \num{3.8e12} +
\num{4.6e12}$ geeft $u = \num{8.4e12}\times\num{1.5e5}/(\num{6e6}\times10^{20})
= \qty{2.1e-9}{m.s^{-1}} = \qty{6.6}{cm/yr}$, de snelheid van de snelle oceanische
platen.
\end{solution}

\begin{solution}{pb:b3:mantle-convection:1}
\textbf{1.} $p = 1000\times9.8\times300 = \qty{2.94}{MPa}$.

\textbf{2.} $w = \rho_wh/\rho_m = 300\times1000/3300 = \qty{90.9}{m}$. De korst
zakt als een blok omlaag; wat onder de belasting wegstroomt en weer terugkeert, is de
mantel, dus de kolom wordt gecompenseerd door op diepte verplaatste mantel.

\textbf{3.} $\pi\times(\num{1.25e5})^2 = \qty{4.91e10}{m^2}$, en het water
weegt $\num{4.91e10}\times300\times1000 = \qty{1.47e16}{kg}$.

\textbf{4.} De kustlijn vormde zich horizontaal, toen de bodem het diepst was ingedrukt in
het midden; sindsdien is het midden meer omhooggekomen dan de rand, zodat de lijn nu
naar het midden oploopt, met \qty{77.2}{m} over \qty{125}{km}:
\qty{0.62}{m/km}.

\textbf{5.} $\lambda/2 = \qty{250}{km}$, dus $\lambda = \qty{500}{km}$ en $k =
2\pi/\num{5e5} = \qty{1.257e-5}{m^{-1}}$.

\textbf{6.} $1/k = \lambda/2\pi = \qty{79.6}{km}$.

\textbf{7.} $w = 90.9 - \text{opheffing}$: 79.4, 69.3, 52.9, 40.4, 26.9 en
\qty{13.7}{m}.

\textbf{8.} $w = w_0e^{-t/\tau}$ geeft $\ln w = \ln w_0 - t/\tau$, een rechte lijn
met helling $-1/\tau$.

\textbf{9.} $\tau = (14 - 1)/\ln(79.4/13.7) = 13/1.7571 = \qty{7.40}{kyr}$.

\textbf{10.} $w(6) = 79.4\,e^{-5/7.40} = \qty{40.4}{m}$, een opheffing van
\qty{50.5}{m}, zoals in de tabel.

\textbf{11.} \qty{13.7}{m} blijven over; $\dot w = 13.7/7400 = \qty{1.85}{mm/yr}$.

\textbf{12.} De volledige \qty{13.7}{m}, steeds langzamer benaderd.

\textbf{13.} De snelheid daalt als $e^{-t/\tau}$: $t = 7.40\ln(1.85/0.5) =
7.40\times1.309 = \qty{9.7}{kyr}$, ruim 9700 jaar extra.

\textbf{14.} $\tau = 7399\times\num{3.156e7} = \qty{2.335e11}{s}$, en $\eta =
\tau\rho_mg\lambda/(4\pi) = \num{2.335e11}\times3300\times9.8\times\num{5e5}/
12.566 = \qty{3.0e20}{Pa.s}$.

\textbf{15.} Net onder het onderste eind van het Fennoscandische bereik. De stroming
onder een belasting met een golflengte van \qty{500}{km} vervalt over \qty{80}{km}:
het herstel bemonstert de ondiepste, zwakste mantel, onder en deels binnen wat
in de werkelijke Aarde de lithosfeer is.

\textbf{16.} $\tau = 4\pi\times\num{3.0e20}/(3300\times9.8\times\num{3e6}) =
\qty{3.9e10}{s} = \qty{1.2}{kyr}$. Onder een belasting zo breed reikt de stroming
tot \qty{477}{km} en dieper, naar een onderste mantel die tien keer viskeuzer is,
wat haar vertraagt.

\textbf{17.} De lithosfeer draagt een deel van een smalle belasting door buiging: de insinking bij
evenwicht is ondieper en breder dan de isostatische \qty{90.9}{m}, en een deel van
het herstel is elastisch en onmiddellijk. Een uniforme viskeuze halfruimte, aangepast aan
de opheffing geeft dan een viscositeit die beide effecten vermengt.

\textbf{18.} $\tau_M = \num{3.0e20}/\num{7.7e10} = \qty{3.9e9}{s}$, ongeveer
\qty{124}{yr}, zestig keer korter dan $\tau$: over duizenden jaren reageert de
mantel als een viskeuze vloeistof.

\textbf{19.} $1 - e^{-50/7.40} = 0.9988$: \qty{99.9}{\percent} van de
insinking bij evenwicht was bereikt, en de aanname klopt.

\textbf{20.} $\eta \propto \lambda$: halveren van $\lambda$ halveert $\eta$, en
$\log_{10}\eta$ daalt met $\log_{10}2 = 0.30$.

\textbf{21.} $\mathrm{Ra} = 3700\times9.9\times\num{3.1e-5}\times1450\times
(\num{6.7e5})^3/(10^{-6}\times\num{3.0e20}) = \num{1.65e6}$.

\textbf{22.} $\num{1.65e6}/657.5 \approx 2500$ keer de kritische waarde.

\textbf{23.} $(\mathrm{Ra}/2\mathrm{Ra}_c)^{1/3} = (1254)^{1/3} = 10.78$: $\delta
= 670/10.78 = \qty{62}{km}$ en $\mathrm{Nu} = 5.4$.

\textbf{24.} $u = (10^{-6}/\num{6.7e5})\times10.78^2 = \qty{1.73e-10}{m.s^{-1}}
= \qty{0.55}{cm/yr}$, twee tot veertien keer trager dan de platen.

\textbf{25.} $q = 5.4\times3.1\times1450/\num{6.7e5} = \qty{36}{mW.m^{-2}}$. De
laag laat de warmte weg die binnen de mantel wordt geproduceerd en de warmte uit de
kern en de onderste mantel; en de schaling laat de constanten van orde één weg,
terwijl de platen, die als halfruimten vanaf de ruggen afkoelen, op jonge zeebodem
veel meer warmte kwijtraken dan een gemiddelde grenslaag zou doen.

\textbf{26.} $\eta = \qty{3.0e20}{Pa.s}$: $\log_{10}\eta =
\textbf{\num[text-series-to-math=true]{20.48}}$.
\end{solution}
